\documentclass{exos}
\usepackage{main}

\begin{document}
Pour étudier la propagation de rumeur, on peut utiliser un graphe pour illustrer le réseau de diffusion de l'information. Dans le schéma suivant, une fake news va être propagée par $A$. Les nombres correspondent au temps que l'information met à être transmise. Ici, on suppose que quiconque transmet la fake news qu'il vient d'entendre à tous ses voisins en même temps.
\textbf{Au bout de combien de temps $F$ reçoit-il la fake news ?}

\begin{center}
\begin{tikzpicture}[scale=1.5]
\tikzstyle{node}=[circle,draw]
\node[node] (A) at (0,0) {$A$};
\node[node] (B) at (2,-1) {$B$};
\node[node] (C) at (1,1.5) {$C$};
\node[node] (D) at (-1,2) {$D$};
\node[node] (E) at (3,3) {$E$};
\node[node] (F) at (4,1) {$F$};

\draw[dashed] (A) -- (B) node[midway] {$7$};
\draw[dashed] (A) -- (C) node[midway] {$9$};
\draw[dashed] (A) -- (D) node[midway] {$14$};
\draw[dashed] (B) -- (C) node[midway] {$10$};
\draw[dashed] (D) -- (C) node[midway] {$2$};
\draw[dashed] (D) -- (E) node[midway] {$9$};
\draw[dashed] (B) -- (F) node[midway] {$15$};
\draw[dashed] (C) -- (F) node[midway] {$11$};
\draw[dashed] (E) -- (F) node[midway] {$6$};
\end{tikzpicture}
\end{center}

\vspace*{1cm}
\hrulefill
\vspace*{1cm}

Pour étudier la propagation de rumeur, on peut utiliser un graphe pour illustrer le réseau de diffusion de l'information. Dans le schéma suivant, une fake news va être propagée par $A$. Les nombres correspondent au temps que l'information met à être transmise. Ici, on suppose que quiconque transmet la fake news qu'il vient d'entendre à tous ses voisins en même temps.
\textbf{Au bout de combien de temps $F$ reçoit-il la fake news ?}

\begin{center}
\begin{tikzpicture}[scale=1.5]
\tikzstyle{node}=[circle,draw]
\node[node] (A) at (0,0) {$A$};
\node[node] (B) at (2,-1) {$B$};
\node[node] (C) at (1,1.5) {$C$};
\node[node] (D) at (-1,2) {$D$};
\node[node] (E) at (3,3) {$E$};
\node[node] (F) at (4,1) {$F$};

\draw[dashed] (A) -- (B) node[midway] {$7$};
\draw[dashed] (A) -- (C) node[midway] {$9$};
\draw[dashed] (A) -- (D) node[midway] {$14$};
\draw[dashed] (B) -- (C) node[midway] {$10$};
\draw[dashed] (D) -- (C) node[midway] {$2$};
\draw[dashed] (D) -- (E) node[midway] {$9$};
\draw[dashed] (B) -- (F) node[midway] {$15$};
\draw[dashed] (C) -- (F) node[midway] {$11$};
\draw[dashed] (E) -- (F) node[midway] {$6$};
\end{tikzpicture}
\end{center}
\end{document}